\documentclass[UTF8,11pt,a4paper,oneside]{ctexbook}

\usepackage[a4paper,margin=2.35cm]{geometry}
\usepackage{booktabs}
\usepackage{longtable}
\usepackage{array}
\usepackage{tabularx}
\usepackage{enumitem}
\usepackage{xcolor}
\usepackage{hyperref}
\usepackage{fancyhdr}
\usepackage{titlesec}
\usepackage{fvextra}
\usepackage{setspace}
\usepackage{graphicx}
\usepackage{url}

\hypersetup{
  colorlinks=true,
  linkcolor=blue!50!black,
  urlcolor=blue!50!black,
  pdftitle={构建你的 Superproject：面向复杂系统 Bring-up 的方法学白皮书},
  pdfauthor={OpenAI Codex},
  pdfsubject={Superproject methodology whitepaper},
  pdfkeywords={superproject, bring-up, submodule, strict gate, test-driven, methodology}
}

\setmainfont{Times New Roman}
\setsansfont{Helvetica Neue}
\setmonofont{Menlo}
\setCJKmainfont{PingFang SC}
\setCJKsansfont{PingFang SC}
\setCJKmonofont{STHeiti}

\pagestyle{fancy}
\fancyhf{}
\lhead{构建你的 Superproject}
\rhead{\thepage}
\setlength{\headheight}{14pt}

\setlist[itemize]{leftmargin=2em}
\setlist[enumerate]{leftmargin=2.4em}
\setstretch{1.2}

\DefineVerbatimEnvironment{CodeBlock}{Verbatim}{
  fontsize=\small,
  breaklines=true,
  breaksymbolleft={},
  commandchars=\\\{\}
}

\titleformat{\chapter}{\Huge\bfseries}{第\thechapter 章}{1em}{}
\titleformat{\section}{\Large\bfseries}{\thesection}{0.8em}{}
\titleformat{\subsection}{\large\bfseries}{\thesubsection}{0.8em}{}

\begin{document}

\frontmatter

\begin{titlepage}
  \centering
  \vspace*{2.6cm}
  {\Huge\bfseries 构建你的 Superproject\par}
  \vspace{0.8cm}
  {\LARGE 面向复杂系统 Bring-up 的方法学白皮书\par}
  \vspace{1.6cm}
  {\Large 从 LinxISA 案例出发，但服务于任何需要组织多模块、多仓库、多团队系统的人\par}
  \vspace{2.4cm}
  \begin{minipage}{0.82\textwidth}
    \centering
    \large
    这不是一份关于如何学习 LinxISA 细节的文档。\\
    这是一份关于如何学会构建你自己的 superproject，\\
    如何让复杂系统在多人协作、持续演化和高压验证下仍然保持可收敛、可复现、可治理的白皮书。
  \end{minipage}
  \vfill
  {\large 版本：2026-03-14\par}
  {\large 生成方式：LaTeX / Tectonic\par}
\end{titlepage}

\chapter*{摘要}
\addcontentsline{toc}{chapter}{摘要}

大多数复杂系统并不是死于单点技术难题，而是死于组织失控：代码分散在多个仓库，
边界不清；验证结果散落在聊天记录、日志片段和口头共识中；每次“集成”都像一次
没有航图的夜间飞行。系统似乎一直在前进，但没有人真正知道它是否稳定，也没有人
能在两周之后解释它为何曾经稳定。

superproject 的价值，就在于把这种混乱重新组织成秩序。它不是一个更大的仓库，
也不是一个更长的 README；它更像一座机场塔台：真正飞行的是各架飞机，真正承担
空气动力学的是各自机体，但所有飞机的进近顺序、跑道占用、冲突消解、天气通报和
起降证据，必须由塔台统一协调。没有塔台，每架飞机都可能在局部上是优秀的，但整
个机场一定是危险的。

LinxISA 提供了一个很强的案例：编译器、模拟器、Linux、RTL、模型、libc、文档、
测试矩阵和技能系统并行演进，但最终依靠的不是“多沟通”，而是四个控制平面，
即 topology、ownership、validation 与 evidence。这四个平面共同构成了一套
可迁移的方法学：任何想要构建自己的 superproject 的团队，都可以借鉴其骨架。

本文的目标，是教会读者学会这种骨架。读完之后，读者应该能够：

\begin{itemize}
  \item 判断自己是否需要一个真正的 superproject；
  \item 为自己的系统设计目录与仓库拓扑；
  \item 建立模块边界、所有权和 repin 纪律；
  \item 把“测试很多”升级成“严格 gate 管理”；
  \item 让 agent 与 maintainer 在同一套 machine-readable 规则下工作；
  \item 把一个复杂系统从“能跑”带到“能长期 bring up”。
\end{itemize}

\tableofcontents

\mainmatter

\chapter{为什么复杂系统需要 Superproject}

\section{复杂系统真正难的不是写代码}

单个模块的复杂性，往往还能够由模块自己的代码结构、测试框架和维护者经验来承受。
但一旦系统演化到多仓库、多语言、多执行面、多团队协作的阶段，难点会发生转移。
真正的挑战不再是“这一段代码如何实现”，而变成：

\begin{itemize}
  \item 哪个模块应该承担哪个语义；
  \item 当合同变化时，谁必须同步；
  \item 哪些测试能代表局部正确，哪些 gate 才代表系统正确；
  \item 哪次失败是叶子模块的责任，哪次失败是集成层的责任；
  \item 今天的绿色结果，明天是否还能被复现。
\end{itemize}

如果说单模块工程的重点是“造船”，那么 superproject 面对的是“建港”。造船需要
处理材料、结构和动力；建港则必须同时处理航道、泊位、调度、验收、维修和记录。
没有港口治理，再好的船也会互相阻塞。

\section{什么是 Superproject}

superproject 不是“所有代码都在一个仓库里”的同义词。真正的 superproject 至少有
四个特征：

\begin{enumerate}
  \item 它协调多个真实独立的实现单元；
  \item 它维护这些单元之间的合同、版本关系与验证关系；
  \item 它拥有比单个模块更高一级的 gate 和证据模型；
  \item 它能够回答“整个系统在什么意义上是绿的”。
\end{enumerate}

在这个定义下，superproject 更像一部宪法，而不是一部法典全集。它不亲自规定每个
模块内部如何写每一行代码，但它规定边界、职责、冲突解决方式和通过标准。

\section{如果没有它，会发生什么}

当系统没有 superproject 方法学时，团队通常会落入三种幻觉：

\begin{itemize}
  \item \textbf{局部绿色幻觉}：每个模块都能给出自己的“通过”截图，但没有人能证明整
  个系统一起通过；
  \item \textbf{近期可用幻觉}：今天的工作树可以运行，但一周后无法复现，也不知道究竟
  是谁改变了什么；
  \item \textbf{责任模糊幻觉}：一旦出现回归，所有人都觉得自己“可能没问题”，最后没
  人真正负责系统收敛。
\end{itemize}

这三种幻觉，正是 superproject 要消除的对象。

\chapter{从 LinxISA 案例提炼出的四个控制平面}

\section{为什么不是一张总表，而是四个平面}

很多项目会试图用一张总表管理一切：模块列表、测试结果、负责人、进展备注、风险
提示都混在一起。这样的表在最初看起来很方便，但当系统规模增长后，它会变成一块
巨大的湿泥地：什么信息都在上面，却没有任何一种信息能被严格验证。

LinxISA 的经验表明，更有效的做法是把治理拆成四个平面。它们像建筑中的四个不同
层次：地基、承重结构、安检系统和监控系统。你不能指望用监控摄像头去替代承重墙。

\section{Topology：先决定道路，不要先讨论交通}

Topology 平面关注路径、目录、仓库链接和 canonical location。它回答的是：
哪些地方允许存在，哪些地方禁止存在，新的工作应落到哪里。

这是方法学里最容易被轻视的一层，因为它看起来不像“高技术”。但事实上，topology
是复杂系统的地形学。没有明确地形，所有后续治理都会漂浮。

LinxISA 的做法很有代表性：

\begin{itemize}
  \item 用根目录 \path{.gitmodules} 统一维护 submodule 拓扑；
  \item 用 \path{AGENTS.md} 和 \path{docs/project/navigation.md} 声明允许路径与
  禁止路径；
  \item 用 \path{tools/ci/check_repo_layout.sh} 把目录约束变成 CI 能执行的规则。
\end{itemize}

可迁移原则是：\textbf{先定义地形，再允许建设。} 如果一个团队在 topology 上是
放任的，那么再严格的 test discipline 也只能在流沙上行走。

\section{Ownership：没有明确负责者，系统就没有真正的记忆}

Ownership 平面不是“谁最懂某个模块”，而是“当某种失败发生时，系统知道应该找谁”。
它要求责任边界能被机器读取，而不是只存在于人脑中。

在 LinxISA 里，这一层通过 \path{manifest.yaml}、checklists 和 waiver ledger
来实现。它的启发是：\textbf{所有权必须结构化，才能规模化。}

如果你的系统没有 ownership 平面，团队就像一支没有站位表的足球队：每个人都在跑，
但没有人知道哪个球是自己必须接的。

\section{Validation：把“我们觉得差不多了”升级为严格 gate}

Validation 平面是 superproject 的安检系统。它不问“代码多不多”，也不问“大家辛
不辛苦”，只问一件事：系统是否满足进入下一阶段的客观条件。

LinxISA 的验证层不是一个单一脚本，而是一组彼此配合的脚本：

\begin{itemize}
  \item 结构性静态检查；
  \item AVS 合同检查；
  \item PR-tier 与 nightly-tier strict closure；
  \item dual-lane runtime convergence；
  \item gate consistency 与 freshness 审计。
\end{itemize}

这说明一个重要原则：\textbf{复杂系统的验证不是一把尺子，而是一套闸门。}
就像水库闸门一样，它们不是为了“限制你前进”，而是为了保证系统在更高水位下不被
冲垮。

\section{Evidence：如果没有证据包，绿色就只是情绪}

Evidence 平面让系统具备可追溯性。它要求每次重要 run 都留下足够完整的证据：
命令、lane、run-id、SHA manifest、log 路径、summary JSON，以及与状态页一致
的 canonical report。

LinxISA 的经验在这里尤其值得借鉴：它明确把 markdown 状态页定义为“派生视图”，
而不是 source of truth。这一选择很重要，因为它切断了“手动修 status page”
这种危险捷径。

你可以把 evidence 看成系统的黑匣子。黑匣子不是为了日常美观，而是为了在出现失
败时，系统仍然有记忆。

\chapter{从案例到通用方法：如何构建你自己的 Superproject}

\section{第一步：画出模块地图，而不是先写总脚本}

很多团队上来就想写一个 \path{run_all.sh}，仿佛只要能一键执行，一切就会自然
有序。实际上，一键执行只是最后的皮肤；真正的秩序来自于你先知道系统有哪些模块、
边界在哪里、哪些关系是硬约束、哪些关系只是调用关系。

建议用一张最小模块地图回答下面问题：

\begin{enumerate}
  \item 这个系统有哪些一等模块；
  \item 哪些模块是实现模块，哪些模块是协调模块；
  \item 哪些模块的变化会引发跨模块验证；
  \item 哪些目录或仓库是 canonical location；
  \item 哪些历史路径必须禁止复活。
\end{enumerate}

\section{第二步：为每个模块定义 closure surface}

一个健康的 superproject 不会要求每个模块做“所有测试”，而是要求每个模块提供
可解释的 closure surface。所谓 closure surface，就是“什么样的证据足以证明这
个模块在当前阶段是健康的”。

例如：

\begin{itemize}
  \item 编译器模块可以是 compile suites、coverage 和工具链闭包；
  \item 模拟器模块可以是 runtime suites、strict system 和 opcode sync；
  \item kernel 模块可以是 build closure、smoke boot 和 full boot；
  \item 文档/架构模块可以是 contract lint 和 matrix closure。
\end{itemize}

如果模块没有 closure surface，它就无法被稳定地集成；它只能以“作者自己觉得没问
题”的形式进入系统。这是大多数 superproject 失控的起点。

\section{第三步：把所有权写成机器能懂的格式}

很多团队会说“这个模块归 Alice 管，那个模块归 Bob 管”。这不够。真正需要的是：

\begin{itemize}
  \item 哪个 gate 属于谁；
  \item 哪个 checklist 代表其责任完成；
  \item 哪些 agent 允许跨模块；
  \item 哪类失败可以被 waiver，谁批准，何时过期。
\end{itemize}

一旦写成机器能读取的结构，维护者和 agent 就能在同一套语义下工作。否则，
superproject 规模一增长，就只能靠资深工程师脑中的“暗知识”维持。

\section{第四步：建立双层验证，而不是把所有检查塞进一个脚本}

最可迁移的设计之一，是把验证分成静态层和运行时层。

\subsection{静态层}

静态层关注治理结构本身是否成立。例如：

\begin{itemize}
  \item manifest 是否合法；
  \item gate 是否都有 owner；
  \item checklist ID 是否冲突；
  \item waivers 是否带过期时间；
  \item 模块 scope 是否超界。
\end{itemize}

\subsection{运行时层}

运行时层关注某个具体 run 是否满足通过条件。例如：

\begin{itemize}
  \item 本次 run 的 required row 是否全绿；
  \item lane 和 run-id 是否完整；
  \item summary JSON 是否存在；
  \item 证据是否新鲜且路径正确。
\end{itemize}

把这两层混在一起，会导致一个典型问题：测试失败时，你根本不知道失败的是“系统
行为”还是“治理结构”。分层之后，诊断会清晰很多。

\section{第五步：设计 Repin Discipline}

repin 是 superproject 最危险、也最能暴露成熟度的动作。它看起来只是移动几个 SHA，
但实际上等于宣布“这些独立模块现在可以共同作为一个系统被信任”。

因此，repin 需要纪律：

\begin{enumerate}
  \item 叶子模块先修复并先证明；
  \item 再在 superproject 中更新 submodule SHA；
  \item 再执行严格 cross-repo closure；
  \item 最后才允许合并。
\end{enumerate}

简而言之：\textbf{repin 不是同步动作，而是发布动作。}

\section{第六步：为系统设计证据包}

想让系统具备长期可维护性，你必须让每一次重要结果都能被重新讲述。方法不是要求
大家写更长的总结，而是要求每次 run 都生成证据包。

一个最小证据包应包括：

\begin{itemize}
  \item run-id；
  \item lane；
  \item 精确命令；
  \item SHA manifest；
  \item canonical JSON 中的结果行；
  \item 对应的日志和 summary artifact。
\end{itemize}

一旦有了这种结构，沟通成本会下降很多，因为很多争论都可以回到证据而不是回到记忆。

\chapter{Strict Gate、Strict Check 与 Test-Driven Bring-up 的通用原则}

\section{测试很多，不等于严格}

很多项目测试非常多，但仍然不严格。原因在于“测试集合”与“通过标准”之间缺了一
层治理。strict gate 的本质不是增加测试数量，而是把通过条件形式化。

这意味着：

\begin{itemize}
  \item 哪些 gate 是 required；
  \item 哪些 gate 在哪个 phase 才算 blocker；
  \item waiver 何时合法；
  \item 哪些结果算 stale；
  \item 什么样的 lane policy 才允许宣称系统收敛。
\end{itemize}

\section{Strict Check 是给流程本身加护栏}

strict check 的核心价值，是把“流程是否被正确执行”也变成验证对象。这样一来，
系统不会因为“大家都以为别人会补证据”而出现治理漏洞。

从方法学角度看，strict check 有三个作用：

\begin{enumerate}
  \item 它让流程可审计；
  \item 它让团队不依赖少数资深人员的记忆；
  \item 它让 agent 能在清晰边界内自主执行。
\end{enumerate}

\section{Test-Driven Bring-up 的关键不是先写测试，而是先定义失败面}

很多人把 test-driven 理解成“先写测试再写实现”。在 bring-up 场景里，更准确的
表达是：先定义失败面，再组织修复闭环。

一个成熟的 bring-up loop 应该是：

\begin{enumerate}
  \item 从失败 contract、audit、gate 或 regression 开始；
  \item 用最小命令复现；
  \item 在正确的 owning module 内修复；
  \item 用该模块的 closure surface 验证修复；
  \item 再执行 superproject 集成验证。
\end{enumerate}

如果没有这个 loop，项目就会不断累积“看起来应该没问题”的改动，最终在集成时
一起爆炸。

\chapter{如何真正学会这套方法学}

\section{先学会看边界，再学会看工具}

很多读者在学习 superproject 方法学时，第一反应是去记命令和脚本名字。这当然有
用，但那只是表层。真正应该先学会的是边界感。

你需要训练自己反复问四个问题：

\begin{enumerate}
  \item 这个问题属于哪个模块；
  \item 这个模块的 closure surface 是什么；
  \item 这个失败会不会影响 cross-repo 集成；
  \item 这次结果的权威证据在哪里。
\end{enumerate}

一旦这四个问题成为你的直觉，工具就不再是负担，而会自然落位。

\section{把系统想象成一座城市}

如果你更习惯通过比喻学习，可以把 superproject 想象成一座城市：

\begin{itemize}
  \item 模块是城市里的功能区；
  \item topology 是道路和规划红线；
  \item ownership 是各区的行政与维护责任；
  \item validation 是安检、消防和准入制度；
  \item evidence 是城市的监控、档案与事故记录；
  \item repin 则像修改城市公共基础设施的生效版本。
\end{itemize}

在这样的城市里，真正危险的不是一栋楼有瑕疵，而是消防、道路、产权和档案同时失
效。superproject 的方法学，就是防止这类系统性塌方。

\section{用三层练习法掌握它}

\subsection{第一层：读图}

先不急着改系统。先能读懂：

\begin{itemize}
  \item 模块地图；
  \item gate ownership；
  \item lane 设计；
  \item artifact 布局；
  \item waiver 规则。
\end{itemize}

\subsection{第二层：跟跑}

选择一条失败的 gate，跟着现有流程走完整一遍：

\begin{enumerate}
  \item 复现；
  \item 定位 owning module；
  \item 修复；
  \item 重跑；
  \item 留证据；
  \item 判断是否 repin。
\end{enumerate}

\subsection{第三层：造骨架}

等你能稳定完成第二层，再尝试为自己的项目亲手建立这些文件：

\begin{itemize}
  \item navigation contract；
  \item module manifest；
  \item checklists；
  \item waiver ledger；
  \item canonical report schema；
  \item PR/nightly gate entrypoint；
  \item maintainer repin checklist；
  \item agent SOP。
\end{itemize}

真正学会方法学，不是把别人的脚本背下来，而是能为自己的系统造出同等强度的骨架。

\chapter{给读者的落地蓝图：如何为你的项目搭骨架}

\section{最小可行 Superproject 套件}

如果你要为自己的项目从零开始搭建，建议最少准备八类文件或入口：

\begin{enumerate}
  \item \textbf{导航合同}：声明允许路径、禁止路径、canonical location；
  \item \textbf{模块清单}：列出模块、职责、边界；
  \item \textbf{ownership manifest}：列出 gate 与 owner 的映射；
  \item \textbf{checklists}：给每个域定义可审计的完成标准；
  \item \textbf{waiver ledger}：记录何时允许带着风险前进；
  \item \textbf{strict entrypoint}：PR-tier 和 nightly-tier 入口脚本；
  \item \textbf{canonical report}：机器可读的 gate 报告；
  \item \textbf{agent / maintainer runbook}：给执行者清晰 SOP。
\end{enumerate}

\section{建议的目录雏形}

下面是一种可迁移的最小目录骨架：

\begin{CodeBlock}
superproject/
  modules/
  docs/project/
    navigation.md
    repository-flow.md
    ownership-manifest.yaml
    waivers.yaml
    new-agent-sop.md
    maintainer-repin-checklist.md
  docs/gates/
    latest.json
    logs/
  tools/ci/
    check_repo_layout.sh
  tools/bringup/
    check_static_policy.py
    check_runtime_closure.py
    check_gate_consistency.py
  tools/regression/
    strict_cross_repo.sh
\end{CodeBlock}

\section{建议的最小字段}

ownership manifest 至少应包含：

\begin{itemize}
  \item 模块名；
  \item owner；
  \item 允许修改的路径；
  \item 对应的 checklist；
  \item gate key 与其归属；
  \item 哪些角色允许跨模块。
\end{itemize}

waiver ledger 至少应包含：

\begin{itemize}
  \item gate key；
  \item owner；
  \item issue 或原因；
  \item phase；
  \item \path{expires_utc}；
  \item 是否仍 active。
\end{itemize}

\section{两条最重要的落地建议}

\subsection{不要一开始追求“大而全”}

真正有效的 superproject 往往从很少的硬约束开始，例如：

\begin{itemize}
  \item 禁止路径漂移；
  \item 每个 gate 必须有 owner；
  \item canonical JSON 是最高真相；
  \item repin 必须在 leaf proof 之后。
\end{itemize}

先把这四条做硬，再逐步增加复杂度。

\subsection{不要让流程只存在于“资深工程师知道”的层面}

如果一个系统的关键治理规则无法被新成员或 agent 从仓库中直接读到，这个系统迟早
会在规模增长后失去稳定性。方法学真正成熟的标志，不是“老板觉得放心”，而是
“新人能依靠文档和脚本走到正确位置”。

\chapter{新 Agent 的通用 SOP 模板}

\section{使命定义}

一个新 agent 的第一职责不是“尽快做出修改”，而是先对系统秩序负责。具体来说，
它应做到：

\begin{enumerate}
  \item 保持在 canonical workspace 与声明过的 module scope 内；
  \item 找到最小失败面；
  \item 优先修 leaf module；
  \item 再考虑 repin 与 strict closure；
  \item 留下可复核证据。
\end{enumerate}

\section{进入系统时的第一组动作}

\begin{CodeBlock}
git submodule sync --recursive
git submodule update --init --recursive
bash tools/ci/check_repo_layout.sh
\end{CodeBlock}

如果 layout 不通过，后续运行结果一律先打折扣。

\section{新 Agent 的五条戒律}

\begin{itemize}
  \item 不在未声明模块中越界修改；
  \item 不在 superproject 层掩盖叶子模块缺陷；
  \item 不手改派生状态页；
  \item 不把无 waiver 的红 gate 当作“暂时可接受”；
  \item 不在证据缺失时宣称闭环完成。
\end{itemize}

\section{推荐的标准闭环}

\begin{enumerate}
  \item 读取导航合同、manifest 与 checklist；
  \item 用最小命令复现失败；
  \item 路由到 owning module；
  \item 修复并重跑 leaf gate；
  \item 如需要，repin；
  \item 运行 strict cross-repo closure；
  \item 检查 artifact 与 canonical report。
\end{enumerate}

\chapter{Maintainer 的通用 Repin Checklist 模板}

\section{Repin 前先问自己}

在移动任何 SHA 之前，maintainer 应先问：

\begin{enumerate}
  \item 这次 repin 的业务或技术原因是什么；
  \item 对应 leaf proof 是什么；
  \item 哪些 cross-repo gate 可能受影响；
  \item 这是一笔“明确的发布动作”，还是只是“顺手同步”。
\end{enumerate}

如果答案偏向后者，应该停下来压缩范围。

\section{Maintainer 的八项合并前核对}

\begin{enumerate}
  \item intended modules 是本批次唯一发生 SHA 变化的模块；
  \item 每个模块都有 leaf proof；
  \item PR-tier 或 nightly-tier strict closure 已执行；
  \item 需要时完成 dual-lane 验证；
  \item canonical report 是新鲜的；
  \item 派生状态页与 canonical report 一致；
  \item waiver 已审查；
  \item commit message 能解释 repin 原因与范围。
\end{enumerate}

\section{一句最重要的提醒}

maintainer 不应把 repin 当作“仓库同步”。repin 在方法学上更接近一次小型发布：
它向系统宣布一组版本关系开始生效，因此必须带着证据和纪律进入主线。

\chapter{常见反模式}

\section{路径漂移}

项目越做越久，大家越容易为了“先做成”而在仓库边缘长出许多近似目录。短期看很灵活，
长期看等于让 topology 平面失效。

\section{状态页先行}

一些团队会先改进度页、再补证据，仿佛文字能替代事实。这种做法的危险之处在于，
它会把整个系统训练成“先讲故事，再找证据”。

\section{把多模块一起顺手升级}

这是最典型的 repin 失控模式。它让任何失败都变得难以归因，也让任何成功都失去可
复现性。

\section{把资深工程师当作方法学本身}

如果一位关键工程师请假，系统就几乎无法完成闭环，那么你拥有的不是成熟的方法学，
而只是一个高负载的人类路由器。

\chapter{结论：如何判断你的 Superproject 已经成熟}

你可以用下面五个问题来审视自己的系统：

\begin{enumerate}
  \item 新成员是否能在不依赖口头传承的前提下找到正确路径；
  \item 每个关键 gate 是否都有明确 owner；
  \item 每次重要 run 是否都有完整证据包；
  \item repin 是否始终建立在 leaf proof 之上；
  \item 系统是否能说清“整体绿色”的定义。
\end{enumerate}

如果这五个问题中大部分答案仍然含糊，那么说明你可能拥有多个仓库，但还没有真正
拥有一个 superproject。

而一旦这些问题都能被清楚回答，你就不只是拥有了一套 bring-up 流程，而是拥有了
一套能在复杂度增长中继续生存的方法学。这正是 superproject 的意义：它不是为了
让系统看起来更宏大，而是为了让系统在规模变大之后仍然不失控。

\backmatter

\chapter*{附录：可直接借用的最小命令集合}
\addcontentsline{toc}{chapter}{附录：可直接借用的最小命令集合}

\begin{CodeBlock}
# 1. 初始化与 layout 检查
git submodule sync --recursive
git submodule update --init --recursive
bash tools/ci/check_repo_layout.sh

# 2. 静态策略检查
python3 tools/bringup/check_multi_agent_gates.py \
  --strict-always \
  --mode static \
  --manifest docs/bringup/agent_runs/manifest.yaml \
  --waivers docs/bringup/agent_runs/waivers.yaml \
  --checklists-root docs/bringup/agent_runs/checklists

# 3. 合同检查
python3 tools/bringup/check_avs_contract.py --matrix avs/linx_avs_v1_test_matrix.yaml
python3 tools/bringup/check_avs_profile_closure.py \
  --matrix avs/linx_avs_v1_test_matrix.yaml \
  --status avs/linx_avs_v1_test_matrix_status.json \
  --tier pr

# 4. PR-tier strict closure
LINX_GATE_TIER=pr RUN_EXTENDED_CROSS_GATES=1 \
bash tools/regression/strict_cross_repo.sh

# 5. Dual-lane convergence
LINX_GATE_TIER=pr RUN_EXTENDED_CROSS_GATES=1 \
bash tools/bringup/run_runtime_convergence.sh --lane pin --run-id <run-id-pin>
LINX_GATE_TIER=pr RUN_EXTENDED_CROSS_GATES=1 \
bash tools/bringup/run_runtime_convergence.sh --lane external --run-id <run-id-external>
\end{CodeBlock}

\end{document}
